\documentclass[12pt]{article}
\usepackage[utf8]{inputenc}
\usepackage[T1]{fontenc}
\usepackage{hyperref}
\title{Kaulong Language Resources for K-12 Curriculum Development}
\author{Compiled by Claude Sonnet 5.5 (claude-sonnet-5-5) for Andr\'as Kornai}
\date{}
\begin{document}
\maketitle

\section{Introduction}
Kaulong (ISO 639-3 \texttt{pss}, Glottocode kaul1240, also called Pasismanua) is an Austronesian (Oceanic) language of the south-west of Kandrian District, West New Britain Province, Papua New Guinea, with roughly 4{,}000 speakers according to the sources found. It is not known to be taught in schools. Web search turned up very little: no primer, textbook, dictionary or community programme was found. This dossier therefore lists only what was actually seen.

\section{Primary Educational Resources}
No primers, textbooks, readers or school programmes for Kaulong were found in search results.

\section{Linguistic and Reference Materials}
\begin{itemize}
\item Wikipedia, ``Kaulong language'': overview, phonology (consonant and vowel inventory, diphthongs; citing a 1992 SIL phonology dataset by Craig Throop) and a pointer to a grammar sketch by Lynch, Ross and Crowley in \emph{The Oceanic Languages} (2002). \url{https://en.wikipedia.org/wiki/Kaulong_language}
\item Verbix page with a short list of Kaulong verbs (very limited). \url{https://docs.verbix.com/Languages/Kaulong}
\item Glosbe crowd-sourced translation pairs involving Kaulong (not a formal dictionary; quality unverified). \url{https://glosbe.com/aud/pss}
\end{itemize}

\section{Religious and Community Resources}
None found.

\section{Digital and Audio Resources}
No audio, video or corpus specific to Kaulong was confirmed. A PARADISEC catalogue item on New Britain (\url{https://catalog.paradisec.org.au/collections/PAMBU/items/DOC615}) appeared in results, but it was not tied to Kaulong and is not claimed as a Kaulong resource.

\section{Curriculum Development Context}
Materials are extremely scarce. Any curriculum would need new documentation work with the community, starting from the published grammar sketch and phonology data. Checking the \emph{Oceanic Languages} volume and SIL/Throop materials directly is the suggested next step.

\end{document}